\chapter{Closed Systems and Memory: Finite Populations, MVA, Campbell's Theorem, Eviction}\label{sec:L4}

All models so far were \term{open systems}. Customers arrive from an infinite
outside population, leave after service, and the arrival rate does not change
when the queue is long. Agent sessions are different. A session submits a turn,
which passes the prefill instance, the link and a decode instance, runs a tool on
the response, and only then submits its next turn. A
waiting session cannot create a new turn, so the longer the queue, the fewer the
arrivals. This lecture treats this \emph{closed} structure
(\cref{L4:sec:closed}), computes the mean memory held by sessions
(\cref{L4:sec:campbell}), and solves the combinatorial optimisation problem of
what to drop when memory runs short (\cref{L4:sec:knapsack}).

\section{Closed systems and mean value analysis}\label{L4:sec:closed}

\begin{definition}[Delay station, think time]\label{L4:def:delay}
A \term{delay station} or \term{infinite-server station} is a station with
infinitely many servers, so that an arriving customer is served at once without
waiting. A customer's sojourn time equals its own service time and does not
depend on other customers. In a closed system the time spent at a delay station
is called the \term{think time}, and its mean is written $Z$.
\end{definition}

\begin{definition}[Closed system, finite-source model]\label{L4:def:closed}
A \term{closed system} is a system in which a fixed number $N$ of customers
circulate among stations forever, with no entry from or exit to the outside. The
\term{finite-source model} or machine-repairman model M/M/1//$N$ is the closed
system made of two stations:
\begin{enumerate}[label=(\alph*)]
\item a delay station with exponential think times of mean $Z$,
\item a FIFO single server with exponential service times of mean $1/\mu$.
\end{enumerate}
Each customer (session) submits a turn to the server when it finishes thinking,
and starts thinking again when its service is done. All think times and service
times are mutually independent.
\end{definition}

\serving{Sessions are the customers of the closed loop prefill instance $\to$ link
$\to$ decode instance $\to$ tool $\to$ prefill instance. Seen from the prefill
instance, the server is the prefill server and everything else (transfer, decode,
tool call and user response) is the ``think time''. A session waiting for prefill
cannot submit its next turn. A memory pool holds a few million tokens and agent
contexts reach hundreds of thousands of tokens, so an instance serves only about
ten live sessions at a time; at this scale the $N=\infty$ approximation can be badly
wrong.}

Let $X(t)$ be the number of turns at the server (waiting or in service) at time
$t$. There are $N-X(t)$ thinking sessions, each finishing its think time
independently at rate $1/Z$ (\cref{prop:memoryless}), so $X$ is a birth--death
process on $\{0,1,\dots,N\}$ with birth rates $\lambda_n=(N-n)/Z$ and death
rates $\mu_n=\mu$ ($n\ge1$). By \cref{thm:birthdeath} the stationary law is
\begin{equation}\label{L4:eq:piN}
  \pi_N(n)=\frac{1}{G_N}\prod_{k=0}^{n-1}\frac{(N-k)/Z}{\mu}
  =\frac{1}{G_N}\,\frac{N!}{(N-n)!}\Bigl(\frac{1}{\mu Z}\Bigr)^{n},
  \qquad n=0,1,\dots,N,
\end{equation}
with normalising constant $G_N$. The state space is finite, so no stability
condition is needed.

In an open system PASTA (\cref{thm:pasta}) told us that an arriving customer sees
the time-average law. In a closed system arrivals are not Poisson, and an
arriving customer does not see the time-average law. Instead the following
holds.

\begin{theorem}[Arrival theorem \cite{reiser1980,lazowska1984}]\label{thm:arrival}
In M/M/1//$N$ in steady state, the law of the number of turns at the server seen
by an arriving turn (excluding itself), that is, the long-run fraction $a_N(n)$
of arrivals that occur in state $n$, equals the time-stationary law of the same
system with $N-1$ sessions: $a_N(n)=\pi_{N-1}(n)$, $n=0,\dots,N-1$.
\end{theorem}

\begin{proof}
In state $n$ arrivals occur at rate $(N-n)/Z$. By the ergodic theorem
(\cref{L2:thm:ergodic}) and the law of large numbers for jump counts, the number of
arrivals in state $n$ during $[0,t]$, divided by $t$, converges almost surely to
$\pi_N(n)(N-n)/Z$. Hence
\[
  a_N(n)=\frac{(N-n)\,\pi_N(n)}{\sum_{m=0}^{N}(N-m)\,\pi_N(m)} .
\]
By \eqref{L4:eq:piN}, for $n\le N-1$ we have
$(N-n)\frac{N!}{(N-n)!}=N\frac{(N-1)!}{(N-1-n)!}$, and for $n=N$,
$(N-n)\pi_N(n)=0$. Therefore
$a_N(n)\propto\frac{(N-1)!}{(N-1-n)!}(\mu Z)^{-n}\propto\pi_{N-1}(n)$, and since
both are probability distributions they are equal.
\end{proof}

\begin{remark}[Back to PASTA as $N\to\infty$]\label{L4:rem:arrival-limit}
Fix $\lambda>0$ and $\mu>\lambda$, and let the think time grow with the population,
$Z=N/\lambda$, so that the total arrival rate when all sessions think stays $\lambda$.
Then $\pi_N(n)\propto\frac{N!}{(N-n)!\,N^n}\rho^n$ with $\rho=\lambda/\mu$, and
$\frac{N!}{(N-n)!\,N^n}=\prod_{k<n}(1-k/N)$ lies in $[0,1]$ and tends to $1$. By
dominated convergence (with dominating sequence $\rho^n$) both $\pi_N$ and
$\pi_{N-1}$ converge to the M/M/1 law $(1-\rho)\rho^n$. So in the limit the law seen
by an arriving turn equals the time-average law: this is PASTA. For a finite $N$ the
arrival theorem measures exactly how far the closed system is from that limit.
\end{remark}

The arrival theorem lets us solve a closed system by recursion on the number of
sessions. With $n$ sessions, let $Q_n$ be the stationary mean number of turns at
the server, $R_n$ the mean response time of a turn (wait plus service), and $X_n$
the throughput (turns completed per unit time).

\begin{theorem}[Mean value analysis \cite{reiser1980}]\label{thm:mva}
In M/M/1//$n$, $Q_0=0$ and for $n\ge1$
\[
  R_n=\frac{1+Q_{n-1}}{\mu},\qquad X_n=\frac{n}{Z+R_n},\qquad Q_n=X_nR_n .
\]
Hence $Q_n=n(1+Q_{n-1})/(\mu Z+1+Q_{n-1})$, the server utilisation is
$\rho_n=X_n/\mu=n/(\mu Z+1+Q_{n-1})$, and the mean waiting time of a turn
(excluding service) is $W_n=Q_{n-1}/\mu$.
\end{theorem}

\begin{proof}
By \cref{thm:arrival} an arriving turn sees on average $Q_{n-1}$ turns ahead of
it. Under FIFO all of them must finish before its service begins. By
memorylessness (\cref{prop:memoryless}) the remaining time of the turn in service
again has mean $1/\mu$, and each waiting turn has mean $1/\mu$, so the mean wait
is $Q_{n-1}/\mu$, and the response time is this plus the turn's own service
$1/\mu$, which is $R_n$. One cycle of a session takes $Z+R_n$ on average.
Applying Little's law (\cref{thm:little}) to the whole system (all $n$ sessions,
always $n$ customers) gives $n=X_n(Z+R_n)$, and applying it to the server gives
$Q_n=X_nR_n$. The utilisation is $\rho_n=X_n/\mu$ by \cref{cor:busy}, and
substitution gives the remaining formulas.
\end{proof}

\begin{proposition}[Finite population]\label{prop:finite}
Let $c=\mu Z>0$ and define $Q_n$ by the recursion of \cref{thm:mva}.
\begin{enumerate}[label=(\roman*)]
\item $Q_n$ is nondecreasing in $n$, $0\le Q_n\le n$, and $Q_n$ is nonincreasing
  in $c$. Moreover $\rho_N<1$ for all $N\ge1$.
\item $W_N\le\rho_N/(\mu(1-\rho_N))$. The right side is the open M/M/1 waiting
  time at the same utilisation (\cref{prop:mm1}).
\end{enumerate}
\end{proposition}

\begin{proof}
Put $g_c(q)=(1+q)/(c+1+q)$, so that $Q_n=n\,g_c(Q_{n-1})$. On $q\ge0$, $g_c$ is
nondecreasing in $q$, nonincreasing in $c$, and $0<g_c\le1$.

(i) $0\le Q_n\le n$ follows by induction from $g_c\le 1$. $Q_1=g_c(0)\ge0=Q_0$,
and if $Q_n\ge Q_{n-1}$ then $Q_{n+1}=(n+1)g_c(Q_n)\ge n\,g_c(Q_{n-1})=Q_n$.
For $c\le c'$, assuming $Q_{n-1}(c')\le Q_{n-1}(c)$,
$Q_n(c')=n\,g_{c'}(Q_{n-1}(c'))\le n\,g_{c'}(Q_{n-1}(c))\le n\,g_c(Q_{n-1}(c))=Q_n(c)$.
$\rho_N<1$ is equivalent to $Q_{N-1}>N-1-c$, so it suffices to show $Q_n>n-c$ for
all $n\ge0$ by induction. For $n=0$ this is $0>-c$, and a computation gives
\[
  Q_{n+1}-(n+1-c)=\frac{c\,(Q_n+c-n)}{c+1+Q_n}>0 .
\]

(ii) Since $\rho_N<1$, the inequality is equivalent to
$Q_{N-1}(1-\rho_N)\le\rho_N$, that is, $Q_{N-1}\le\rho_N(1+Q_{N-1})=Q_N$, which is
(i).
\end{proof}

\begin{example}\label{L4:ex:two}
With $N=2$ and $Z=2/\mu$ ($c=2$), $Q_1=1/3$, $\rho_2=2/(3+1/3)=3/5$ and
$W_2=1/(3\mu)$. The open formula at the same utilisation gives
$\rho_2/(\mu(1-\rho_2))=3/(2\mu)$, which is $4.5$ times the actual wait.
\end{example}

\begin{remark}
For exponential work and at the same utilisation, the open formula
(\cref{thm:pk}) and the price of a miss built on it (\cref{thm:price})
overestimate the wait (\cref{prop:finite}(ii)). For general work this comparison
has not been proved. Moreover, at most $N$ turns exist, so $0\le L_P\le N$, and no
policy change can raise the mean number of turns at the server by more than
$N-L_P$. This cap does not depend on the work law.
\end{remark}

\begin{remark}
We assumed exponential think times, but the results depend only on their mean
$Z$. The delay station is insensitive like the PS station (\cref{thm:insens}), and
by the BCMP theorem a network of delay stations and PS stations has a
product-form stationary law \cite{bcmp1975}. The service times of a FIFO server
do not generalise in this way.
\end{remark}

\section{M/G/\texorpdfstring{$\infty$}{infinity} and Campbell's theorem}\label{L4:sec:campbell}

We now count memory. Live sessions do not wait for each other and each keeps its
own KV, so from the point of view of memory the sessions are customers of a delay
station. At the prefill instance the KV that stays between turns is the cached
prefix of each paused session; at a decode instance it is the KV of the running
turns.

\begin{definition}[M/G/$\infty$]\label{L4:def:mginf}
An M/G/$\infty$ system is a delay station with Poisson arrivals of rate $\lambda$
(\cref{def:poisson}) and i.i.d.\ service times $S$ independent of the arrival
process.
\end{definition}

The following lemma is a standard property of the Poisson process.

\begin{lemma}[Conditional uniformity \cite[\S2.3]{ross1996}]\label{L4:lem:uniform}
For a Poisson process of rate $\lambda$, the number of arrivals in an interval
$(a,b]$ is Poisson$(\lambda(b-a))$, and conditional on this number being $n$, the
arrival times (as an unordered set) have the law of $n$ i.i.d.\ samples from the
uniform distribution on $(a,b]$.
\end{lemma}

\begin{proposition}[M/G/$\infty$ occupancy \cite{ross1996}]\label{L4:prop:mginf}
If $\E[S]<\infty$, then in an M/G/$\infty$ system started empty at time $-t$, the
number of customers at time $0$ is
Poisson$\bigl(\lambda\int_0^t\Prob(S>u)\dd u\bigr)$, and as $t\to\infty$ it
converges in distribution to Poisson$(\lambda\E[S])$.
\end{proposition}

\begin{proof}
Let $A$ be the number of arrivals in $(-t,0]$. Conditional on $A=n$, the arrival
times are i.i.d.\ uniform (\cref{L4:lem:uniform}) and the service times are
independent of them. A customer who arrived at time $-u$ is still present at time
$0$ only if $S>u$, so each customer remains independently with probability
$p_t=\frac1t\int_0^t\Prob(S>u)\dd u$. Hence the number $X$ remaining satisfies
\[
  \Prob(X=k)=\sum_{n\ge k}e^{-\lambda t}\frac{(\lambda t)^n}{n!}\binom nk p_t^k(1-p_t)^{n-k}
  =e^{-\lambda tp_t}\frac{(\lambda tp_t)^k}{k!} .
\]
Finally $\lambda tp_t=\lambda\int_0^t\Prob(S>u)\dd u\to\lambda\int_0^\infty\Prob(S>u)\dd u=\lambda\E[S]$.
\end{proof}

That the limit depends on the law of $S$ only through its mean is the
insensitivity of the delay station. When the quantity each customer holds (for
example, KV bytes) changes over time, the next theorem gives its mean.

\begin{theorem}[Campbell's mean formula \cite{kingman1993}]\label{thm:campbell}
Let $\{T_k\}$ be the points of a Poisson process of rate $\lambda$ on $\R$, and
attach to each point an i.i.d.\ mark $Y_k\sim Y$ (with values in an arbitrary
measurable space), independent of the points. For measurable $f\ge0$,
\[
  \E\Bigl[\sum_k f(T_k,Y_k)\Bigr]=\lambda\int_{\R}\E[f(t,Y)]\dd t .
\]
\end{theorem}

\begin{proof}
First consider the sum over the points in a bounded interval $I=(a,b]$ only.
Conditional on the number of points being $n$, the pairs $(T_k,Y_k)$ are $n$
i.i.d.\ copies of $(U,Y)$ with $U\sim$ Uniform$(I)$ and $Y$ independent
(\cref{L4:lem:uniform}). Hence
\[
  \E\Bigl[\sum_{T_k\in I}f(T_k,Y_k)\Bigr]
  =\sum_{n\ge0}\Prob(A_I=n)\,n\,\E[f(U,Y)]
  =\lambda(b-a)\cdot\frac{1}{b-a}\int_a^b\E[f(t,Y)]\dd t .
\]
(The second equality uses Tonelli's theorem.) For $I_m=(-m,m]$ both sides are
nondecreasing in $m$, and the monotone convergence theorem gives the claim as
$m\to\infty$.
\end{proof}

\begin{corollary}[Mean resident KV]\label{cor:resident-kv}
Let sessions arrive in a Poisson process of rate $\Lambda$, and let session $k$
have an i.i.d.\ lifetime $T_s^{(k)}$ and KV size $K_k(u)\ge0$ at age $u$ (the
lifetime and the trajectory may be dependent). In steady state (started at time
$-\infty$), the mean total resident KV at a given instant is
\[
  \Lambda\,\E\Bigl[\int_0^{T_s}K(u)\dd u\Bigr].
\]
In particular, if $K(u)\equiv\bar K$ is constant for each session, this equals
$\Lambda\bigl(\E[T_s]\E[\bar K]+\Cov(T_s,\bar K)\bigr)$. It exceeds the naive
product $\Lambda\E[T_s]\E[\bar K]$ if and only if $\Cov(T_s,\bar K)>0$, that is,
when long sessions have large contexts.
\end{corollary}

\begin{proof}
Attach the mark $Y_k=(T_s^{(k)},K_k(\cdot))$ and put
$f(t,(T,K))=K(-t)\1\{0\le -t<T\}$. The resident KV at time $0$ is
$\sum_kf(T_k,Y_k)$. By \cref{thm:campbell} and Tonelli, its mean is
\[
  \Lambda\int_{-\infty}^0\E[K(-t)\1\{-t<T_s\}]\dd t=\Lambda\E\Bigl[\int_0^{T_s}K(u)\dd u\Bigr].
\]
If $K\equiv\bar K$, the integral is $T_s\bar K$, and
$\E[T_s\bar K]=\E[T_s]\E[\bar K]+\Cov(T_s,\bar K)$.
\end{proof}

Taking $K\equiv1$ gives the mean number of sessions $\Lambda\E[T_s]$, in agreement
with Little's law.

\section{Eviction: the covering knapsack problem}\label{L4:sec:knapsack}

When the prefill instance's pool runs short, cached prefixes of paused programs must be
dropped (\term{eviction}); a decode instance that keeps paused KV faces the same problem
with its own prices (\cref{sec:L5}). Program $i$ occupies $c_i>0$ bytes, and dropping it
costs $w_i\ge0$ in expectation (e.g.\ $p_i$ times the price of a miss, \cref{thm:price}).

\begin{definition}[Covering knapsack problem]\label{L4:def:knapsack}
For a finite set $I$ and a target $\Delta C>0$, the \term{covering knapsack
problem} is
\[
  \min_{S\subseteq I}\ w(S):=\sum_{i\in S}w_i\quad\text{s.t.}\quad c(S):=\sum_{i\in S}c_i\ge\Delta C .
\]
Allowing $x\in[0,1]^I$ in place of $S$ and minimising $\sum w_ix_i$ subject to
$\sum c_ix_i\ge\Delta C$ gives the \term{fractional version}, close to what engines
solve, since they drop fixed-size blocks; $v_i=w_i/c_i$ is the \term{price per byte}.
\end{definition}

\begin{theorem}[Threshold rule \cite{dantzig1957}]\label{thm:threshold}
\begin{enumerate}[label=(\roman*)]
\item For every $\theta\ge0$, the set $T_\theta=\{i:v_i\le\theta\}$ costs no more
  than any $S$ with $c(S)\ge c(T_\theta)$.
\item In the fractional version, the solution that sets $x_i=1$ in increasing
  order of $v_i$ and fills the last item partially to meet the target exactly is
  optimal (provided $c(I)\ge\Delta C$).
\end{enumerate}
\end{theorem}

\begin{proof}
(i) On $S\setminus T_\theta$ we have $w_i>\theta c_i$, and on $T_\theta\setminus S$
we have $w_i\le\theta c_i$, so
\[
  w(S)-w(T_\theta)=\sum_{S\setminus T_\theta}w_i-\sum_{T_\theta\setminus S}w_i
  \ge\theta\bigl(c(S)-c(T_\theta)\bigr)\ge0 .
\]
(ii) follows by the same exchange argument with $\theta$ the price per byte of the
last, partially taken item.
\end{proof}

When programs must be dropped whole, price-per-byte order can be arbitrarily bad,
and the standard repair (a greedy guarded by a guess of the most expensive dropped
item) is within a factor 2 of the optimum \cite{csirik1991}. Real systems evict by a
key fixed in advance rather than by a price, such as shortest context first (SF),
least recently used (LRU) or a time-to-live; such keys do not look at the resume
probability $p_i$.

\begin{proposition}[Price-blind keys]\label{prop:blind}
Consider the covering knapsack problem with $w_i=p_i\Phi_i$.
\begin{enumerate}[label=(\roman*)]
\item For every rule that chooses one of two programs from any information except
  their resume probabilities, and every $R$, there are resume probabilities in
  $(0,1]$ under which either program alone meets the target and the program the rule
  chooses costs more than $R$ times the other: no such rule has a constant
  approximation ratio.
\item If every program resumes and $w_i=c_i^2$ (the quadratic term of the recompute
  work), SF is price-per-byte order and costs at most twice the optimum; the factor 2
  is tight (sizes $K$ and $K+1$ with target $K+1$).
\end{enumerate}
\end{proposition}

\begin{proof}[Proof of (i)]
Take two programs of the same size $c$ that agree in everything the rule observes, and
the target $\Delta C=c$. The rule drops a fixed one, $e$, whatever the resume
probabilities; call the other $o$. With equal prices $\Phi$, $p_e=1$ and
$p_o=1/(|R|+2)$, $w_e=\Phi>R\,w_o$, while the optimum is at most $w_o$.
\end{proof}

Even under (ii), SF is not optimal: with contexts $4,5,6$ and target $6$ it drops
$\{4,5\}$ at cost $41$, while $\{6\}$ costs $36$.

\section{Exercises}

\begin{exercise}
For $N=2$ and $c=\mu Z$, compute $\sum_nn\,\pi_2(n)$ directly from
\eqref{L4:eq:piN} and check that it equals $Q_2$ of \cref{thm:mva}.
\emph{Hint:} $G_2=1+2/c+2/c^2$.
\end{exercise}

\begin{exercise}
Show that the stationary mean number of customers in M/G/$\infty$ is $\lambda\E[S]$
without using \cref{L4:prop:mginf}.
\emph{Hint:} at a delay station $W=\E[S]$, so use Little's law.
\end{exercise}

